% EMNLP 2026 - Compositional Skill Routing for LLM Agents
% Using ACL Rolling Review (ARR) format
\documentclass[11pt]{article}

% ARR / EMNLP 2026 style
\usepackage[hyperref]{acl}  % ACL Rolling Review style
\usepackage{times}
\usepackage{latexsym}
\usepackage{amsmath}
\usepackage{amssymb}
\usepackage{graphicx}
\usepackage{booktabs}
\usepackage{multirow}
\usepackage{xcolor}
\usepackage{enumitem}
\usepackage{url}
\usepackage{algorithm}
\usepackage{algorithmic}
\usepackage{subcaption}

% Custom commands
\newcommand{\system}{\textsc{SkillWeaver}}
\newcommand{\benchmark}{\textsc{CompSkillBench}}
\newcommand{\catrecall}{\text{CatR}}
\newcommand{\chaincat}{\text{Chain}_\text{cat}}

\title{Compositional Skill Routing for LLM Agents:\\Decompose, Retrieve, and Compose}

% Anonymous for review
\author{Anonymous}

\begin{document}
\maketitle

\begin{abstract}
Large language model (LLM) agents increasingly rely on external skills---reusable tool specifications that define capabilities such as code generation, data analysis, or API interaction.
While recent work addresses routing queries to a single best-matching skill, real-world tasks often require \emph{composing} multiple skills into coordinated execution plans.
We formalize this as the \textbf{Compositional Skill Routing} problem: given a complex user query and a large skill library, decompose the query into atomic sub-tasks, retrieve the appropriate skill for each sub-task, and compose an executable plan.
We present \system{}, a three-stage framework consisting of (1)~an LLM-based task decomposer, (2)~a bi-encoder skill retriever with FAISS indexing, and (3)~a compatibility-aware DAG planner.
To support evaluation, we introduce \benchmark{}, a benchmark of 300 compositional queries over 2{,}595 curated skills spanning 17 categories, with ground-truth skill chains and difficulty stratification.
Our experiments reveal that \emph{task decomposition quality is the primary bottleneck}: oracle decomposition achieves 99.5\% skill recall, while the best LLM decomposer reaches only 33.9\% at the category level.
We further find that metadata-only retrieval matches or outperforms body-aware retrieval, and that decomposer choice significantly impacts downstream routing quality.
Our analysis provides actionable insights for building compositional skill routing systems and highlights decomposition as the critical frontier for improvement.
\end{abstract}

% ============================================================
\section{Introduction}
\label{sec:intro}

The agent paradigm for large language models (LLMs) has evolved beyond single-turn generation to encompass tool use, planning, and multi-step task execution \cite{schick2023toolformer,qin2023toolllm,patil2023gorilla}.
A key architectural pattern emerging in modern LLM agents is the use of \emph{skills}: modular, reusable tool specifications that define specific capabilities along with instructions for when and how to invoke them \cite{anthropic2025skills}.
As agent skill libraries grow---with repositories already containing thousands of community-contributed skills---a fundamental routing question arises: \emph{given a user query, which skill(s) should the agent invoke?}

Prior work has largely treated skill routing as a single-skill selection problem, where the goal is to match a query to the single best skill \cite{xu2025skillrouter}.
However, real-world user queries frequently require \emph{multiple} skills working together.
For example, the query ``\textit{Download the dataset, transform it, and create visual reports}'' requires at least three distinct skills: an API client for downloading, a data processor for transformation, and a visualization tool for report generation.

We formalize this as the \textbf{Compositional Skill Routing} problem (Figure~\ref{fig:overview}):
given a complex query $q$ and a skill library $\mathcal{S}$, produce an ordered sequence of skills $[s_1, s_2, \ldots, s_k]$ that collectively accomplish the query, where each $s_i$ handles one atomic sub-task.
This formulation captures the compositional nature of real-world tasks while maintaining the modularity benefits of skill-based architectures.

We present \system{}, a three-stage framework that addresses this problem through:
\begin{enumerate}[nosep,leftmargin=*]
  \item \textbf{Decompose}: An LLM-based task decomposer that breaks complex queries into atomic sub-tasks, each requiring exactly one skill.
  \item \textbf{Retrieve}: A bi-encoder retriever that identifies candidate skills for each sub-task using semantic similarity over skill metadata.
  \item \textbf{Compose}: A compatibility-aware planner that selects the best skill per step while considering inter-skill compatibility, producing an executable DAG.
\end{enumerate}

To evaluate compositional skill routing, we construct \benchmark{}, the first dedicated benchmark for this task.
\benchmark{} contains 300 compositional queries over 2{,}595 curated skills spanning 17 functional categories, with ground-truth skill chains and three difficulty levels.
Skills are sourced from 212 public GitHub repositories and deduplicated to ensure quality.

Our experiments yield several key findings:
\begin{itemize}[nosep,leftmargin=*]
  \item \textbf{Decomposition is the bottleneck}: With oracle (ground-truth) decomposition, our retriever achieves 99.5\% skill Recall@1 and perfect category-level accuracy. LLM decomposition drops this to 33.9\% category Recall@1, identifying decomposition as the critical frontier.
  \item \textbf{Metadata suffices for retrieval}: Metadata-only encoding (name + description) matches or outperforms body-aware encoding that includes the full skill specification, suggesting that concise metadata carries the most discriminative signal.
  \item \textbf{Decomposer choice matters}: Qwen2.5-7B significantly outperforms Llama-3.1-8B and Mistral-7B at task decomposition, with the latter two suffering from over-decomposition that degrades downstream retrieval.
\end{itemize}

\begin{figure*}[t]
\centering
\fbox{\parbox{0.95\textwidth}{\centering\vspace{1.5em}
\small
\textbf{Query:} ``Download the dataset, transform it, and create visual reports'' \\[0.8em]
$\downarrow$ \textbf{Stage 1: Decompose} (LLM) $\downarrow$ \\[0.4em]
$[t_1]$ Download dataset $\;\to\;$ $[t_2]$ Transform data $\;\to\;$ $[t_3]$ Create visual reports \\[0.8em]
$\downarrow$ \textbf{Stage 2: Retrieve} (Bi-Encoder + FAISS) $\downarrow$ \\[0.4em]
$t_1 \to$ \{api-client, http-fetch, ...\} \quad $t_2 \to$ \{csv-parser, etl-pipeline, ...\} \quad $t_3 \to$ \{chart-gen, dashboard, ...\} \\[0.8em]
$\downarrow$ \textbf{Stage 3: Compose} (Compatibility-Aware Planning) $\downarrow$ \\[0.4em]
$s_1$: api-client $\;\xrightarrow{\text{seq}}\;$ $s_2$: csv-parser $\;\xrightarrow{\text{seq}}\;$ $s_3$: chart-gen
\vspace{1.5em}}}
\caption{Overview of \system{}. A complex query is decomposed into atomic sub-tasks, each matched to candidate skills via bi-encoder retrieval, then composed into an executable skill chain using compatibility-aware planning.}
\label{fig:overview}
\end{figure*}

% ============================================================
\section{Related Work}
\label{sec:related}

\paragraph{Tool Selection and Routing.}
Tool selection for LLMs has been studied through API retrieval \cite{patil2023gorilla,qin2023toolllm}, tool documentation matching \cite{hao2024toolkengpt}, and hierarchical routing \cite{xu2025skillrouter}.
SkillRouter \cite{xu2025skillrouter} is closest to our work, introducing a bi-encoder approach for matching queries to individual skills and finding that skill body content is the decisive retrieval signal.
However, all these methods select a \emph{single} tool per query.
Our work extends skill routing to the compositional setting, where multiple skills must be orchestrated.

\paragraph{Task Decomposition.}
LLM-based task decomposition has been explored in planning \cite{huang2022language,wang2023plan}, multi-step reasoning \cite{wei2022chain,zhou2022least}, and agentic frameworks \cite{yao2023react,shinn2023reflexion}.
TaskBench \cite{shen2023taskbench} evaluates tool-augmented LLMs on multi-step tasks but focuses on a fixed set of tools rather than retrieval from a large library.
Our decomposition stage builds on these foundations but is specifically designed for the downstream skill retrieval task.

\paragraph{Retrieval-Augmented Generation.}
Dense retrieval using bi-encoders \cite{karpukhin2020dense,ni2022sentence} has been widely applied to knowledge-intensive NLP tasks.
We adapt this paradigm for skill retrieval, using semantic embeddings of skill metadata to enable efficient nearest-neighbor search over large skill libraries.

% ============================================================
\section{Problem Formulation}
\label{sec:problem}

\paragraph{Skill Library.}
A skill library $\mathcal{S} = \{s_1, \ldots, s_N\}$ contains $N$ skills.
Each skill $s_i$ is a tuple $(n_i, d_i, b_i, C_i)$ where $n_i$ is the name, $d_i$ is a natural language description, $b_i$ is the full specification body (instructions, examples, configuration), and $C_i \subseteq \mathcal{C}$ is a set of functional categories from a taxonomy $\mathcal{C}$.

\paragraph{Compositional Skill Routing.}
Given a complex query $q$ that requires multiple capabilities, the goal is to produce:
\begin{enumerate}[nosep,leftmargin=*]
  \item A decomposition $D(q) = [t_1, \ldots, t_K]$ of $K$ atomic sub-tasks.
  \item A skill assignment $\sigma: [t_1, \ldots, t_K] \to \mathcal{S}^K$ mapping each sub-task to a skill.
  \item An execution plan (DAG) $G = (V, E)$ specifying dependencies between steps.
\end{enumerate}

The compositional routing function is $f: q \to (D, \sigma, G)$ optimizing:
\begin{equation}
\max_{D, \sigma, G} \sum_{k=1}^{K} \text{rel}(t_k, \sigma(t_k)) + \lambda \!\!\!\sum_{(i,j) \in E}\!\!\! \text{compat}(\sigma_i, \sigma_j)
\end{equation}
where $\text{rel}(\cdot)$ measures sub-task--skill relevance and $\text{compat}(\cdot)$ measures inter-skill compatibility.

% ============================================================
\section{Method: \system{}}
\label{sec:method}

\system{} implements compositional skill routing through three cascaded stages (Figure~\ref{fig:overview}).

\subsection{Stage 1: Task Decomposition}

Given a complex query $q$, the task decomposer uses an instruction-tuned LLM to produce an ordered list of atomic sub-tasks:
\begin{equation}
D(q) = \text{LLM}(p_{\text{sys}}, p_{\text{user}}(q)) = [t_1, \ldots, t_K]
\end{equation}
where $p_{\text{sys}}$ is a system prompt that instructs the model to decompose queries into sub-tasks, each requiring exactly one skill.
The output is constrained to a JSON array of strings for reliable parsing.

We apply the model's native chat template (e.g., ChatML for Qwen models) to ensure proper instruction following.
A fallback parser handles cases where the LLM produces numbered lists or other non-JSON formats.

\subsection{Stage 2: Skill Retrieval}

For each sub-task $t_k$, we retrieve the top-$m$ candidate skills from the library using a bi-encoder architecture:
\begin{equation}
\text{cand}(t_k) = \text{top-}m_{s \in \mathcal{S}} \; \cos(E_q(t_k), E_s(s))
\end{equation}
where $E_q$ and $E_s$ are the query and skill encoders (shared parameters), implemented using BGE-large-en-v1.5 \cite{xiao2023bge}.

\paragraph{Skill Representation.}
We investigate two encoding strategies:
\begin{itemize}[nosep,leftmargin=*]
  \item \textbf{Metadata-only}: $\text{repr}(s) = n_s \oplus d_s$ (name and description)
  \item \textbf{Body-aware}: $\text{repr}(s) = n_s \oplus d_s \oplus b_s[:L]$ (with truncated body)
\end{itemize}
where $\oplus$ denotes concatenation and $L=2000$ characters.

Embeddings are $L_2$-normalized and indexed using FAISS \cite{johnson2019faiss} for efficient inner product search.

\subsection{Stage 3: Compatibility-Aware Planning}

Given candidates for each sub-task, the planner selects the final skill per step while considering sequential compatibility:
\begin{multline}
\sigma(t_k) = \arg\max_{s \in \text{cand}(t_k)} \; \alpha \cdot \text{sim}(t_k, s) \\
+ (1-\alpha) \cdot \text{compat}(s_{k-1}, s)
\end{multline}
where $\text{compat}(s_a, s_b)$ combines embedding similarity, category overlap, and I/O pattern heuristics.
The execution DAG assumes sequential dependencies by default, which we find covers the vast majority of compositional queries in practice.

% ============================================================
\section{Benchmark: \benchmark{}}
\label{sec:benchmark}

\subsection{Skill Pool Construction}

We collect skills from 212 public GitHub repositories that adopt the SKILL.md specification format \cite{anthropic2025skills}.
The raw collection yields 11{,}011 skills.
We apply the following curation pipeline:

\begin{enumerate}[nosep,leftmargin=*]
  \item \textbf{Deduplication}: Skills with identical names are merged, reducing to 7{,}681 unique skills.
  \item \textbf{Quality filtering}: Skills without descriptions or with descriptions shorter than 10 characters are removed.
  \item \textbf{Categorization}: Skills are assigned to 17 functional categories (Table~\ref{tab:categories}) using keyword matching on name and description fields only.
  Anti-keywords prevent false positive assignments (e.g., ``debug'' skills are excluded from ``deployment'').
  Skills matching no category or multiple conflicting categories are excluded.
  \item The final pool contains \textbf{2{,}595} categorized skills.
\end{enumerate}

\begin{table}[t]
\centering
\small
\setlength{\tabcolsep}{4pt}
\begin{tabular}{lrl}
\toprule
\textbf{Category} & \textbf{Count} & \textbf{Examples} \\
\midrule
Code Generator & 389 & prd-generator, nextjs-dev \\
API Client & 312 & stripe-api, github-api \\
Test Writer & 248 & unit-test-gen, qa-suite \\
Data Processor & 221 & csv-parser, etl-pipeline \\
Code Analyzer & 187 & code-review, lint-checker \\
Writer & 176 & blog-writer, doc-gen \\
Deployment & 142 & k8s-deploy, ci-cd-setup \\
Security & 118 & vuln-scanner, audit-tool \\
Visualizer & 97 & chart-gen, dashboard-builder \\
Search & 96 & web-search, code-search \\
\midrule
\multicolumn{2}{l}{\textit{+ 7 more categories}} & 609 total \\
\bottomrule
\end{tabular}
\caption{Top 10 skill categories in \benchmark{} (of 17 total). The full pool contains 2{,}595 skills from 212 repositories.}
\label{tab:categories}
\end{table}

\subsection{Query Generation}

Compositional queries are generated by combining skills from different categories into multi-step tasks:

\paragraph{Difficulty Levels.}
\begin{itemize}[nosep,leftmargin=*]
  \item \textbf{Easy} (117 queries): 2 skills, 2 categories
  \item \textbf{Medium} (120 queries): 3 skills, 3 categories
  \item \textbf{Hard} (63 queries): 4--5 skills, 4--5 categories
\end{itemize}

Each query is associated with ground-truth sub-task descriptions (derived from actual skill descriptions), ground-truth skill IDs, required categories, and a sequential execution order.
The benchmark totals 300 queries involving 523 unique ground-truth skills.

\subsection{Evaluation Metrics}

We evaluate at three granularities:

\paragraph{Step-Level Metrics.}
\begin{itemize}[nosep,leftmargin=*]
  \item \textbf{Skill Recall@$k$} (R@$k$): Fraction of steps where the ground-truth skill appears in the top-$k$ candidates.
  \item \textbf{Category Recall@$k$} ($\catrecall$@$k$): Fraction of steps where \emph{any skill from the correct category} appears in the top-$k$.
  This relaxed metric is more practical, as many skills within a category are functionally interchangeable.
\end{itemize}

\paragraph{Chain-Level Metrics.}
\begin{itemize}[nosep,leftmargin=*]
  \item \textbf{Chain Exact Match}: Fraction of queries where \emph{all} steps select the exact ground-truth skill.
  \item \textbf{Chain Category Match} ($\chaincat$): Average fraction of steps per query that select a skill from the correct category.
\end{itemize}

\paragraph{Decomposition Accuracy.} Fraction of queries where the predicted number of sub-tasks matches the ground truth.

% ============================================================
\section{Experimental Setup}
\label{sec:setup}

\paragraph{LLM Decomposers.}
We compare three open-weight 7--8B instruction-tuned models:
Qwen2.5-7B-Instruct \cite{qwen2.5}, 
Llama-3.1-8B-Instruct \cite{dubey2024llama3}, and 
Mistral-7B-Instruct-v0.3 \cite{jiang2023mistral}.
Generation uses temperature $\tau=0.1$, $\text{top}_p=0.9$, and a maximum of 256 new tokens.

\paragraph{Retriever.}
BGE-large-en-v1.5 \cite{xiao2023bge} (1024-dim) serves as the bi-encoder, with FAISS IndexFlatIP for exact inner product search.
We set $k=10$ for retrieval unless otherwise noted.

\paragraph{Baselines.}
\begin{itemize}[nosep,leftmargin=*]
  \item \textbf{Single-Skill}: No decomposition; the full query is used to retrieve one skill set applied to all steps. This represents current single-skill routing methods.
  \item \textbf{Oracle Decomposition}: Ground-truth sub-task descriptions are used for retrieval, establishing an upper bound on decomposition quality.
\end{itemize}

\paragraph{Hardware.} 
All experiments run on a single NVIDIA A100-80GB GPU with 491GB system RAM.

% ============================================================
\section{Results}
\label{sec:results}

\subsection{Main Results}

Table~\ref{tab:main} presents the main experimental results across all configurations.

\begin{table*}[t]
\centering
\small
\resizebox{\textwidth}{!}{
\begin{tabular}{llccccccc}
\toprule
\textbf{Method} & \textbf{Signal} & \textbf{Decomp\%} & \textbf{R@1} & \textbf{R@10} & $\catrecall$\textbf{@1} & $\catrecall$\textbf{@10} & \textbf{Chain}$_E$ & $\chaincat$ \\
\midrule
\multicolumn{9}{l}{\textit{Baselines}} \\
Single-Skill & Meta & -- & 0.000 & 0.029 & 0.211 & 0.698 & 0.000 & 0.216 \\
Single-Skill & Body & -- & 0.002 & 0.028 & 0.285 & 0.633 & 0.000 & 0.293 \\
Oracle Decomp & Meta & 1.000 & \textbf{0.995} & \textbf{1.000} & \textbf{1.000} & \textbf{1.000} & \textbf{0.987} & \textbf{1.000} \\
Oracle Decomp & Body & 1.000 & 0.990 & 1.000 & 0.998 & 1.000 & 0.970 & 0.998 \\
\midrule
\multicolumn{9}{l}{\textit{Full Pipeline (\system{})}} \\
Qwen2.5-7B & Meta & 0.360 & 0.005 & 0.034 & 0.339 & 0.649 & 0.000 & 0.316 \\
Qwen2.5-7B & Body & 0.370 & 0.008 & 0.029 & \underline{0.363} & \underline{0.662} & 0.000 & \underline{0.345} \\
Llama-3.1-8B & Meta & 0.000 & 0.002 & 0.016 & 0.216 & 0.507 & 0.000 & 0.210 \\
Mistral-7B & Meta & 0.177 & 0.005 & 0.036 & 0.243 & 0.607 & 0.000 & 0.248 \\
\bottomrule
\end{tabular}
}
\caption{Main results on \benchmark{}. \textbf{Bold}: best overall. \underline{Underline}: best among full pipeline configurations. R@$k$: Skill Recall@$k$. $\catrecall$@$k$: Category Recall@$k$. Chain$_E$: Chain Exact Match. $\chaincat$: Chain Category Match. Decomp\%: decomposition accuracy (correct number of sub-tasks).}
\label{tab:main}
\end{table*}

\paragraph{Decomposition is the bottleneck.}
The most striking result is the enormous gap between oracle decomposition and LLM decomposition.
Oracle decomposition achieves R@1 of 99.5\% and perfect category accuracy, demonstrating that the retriever is highly effective when given accurate sub-task descriptions.
The best LLM decomposer (Qwen2.5-7B) reaches only $\catrecall$@1 = 36.3\%, indicating that decomposition quality---not retrieval capability---is the primary bottleneck.

\paragraph{Compositional routing outperforms single-skill.}
Despite imperfect decomposition, our full pipeline (Qwen + meta) achieves $\catrecall$@1 = 33.9\%, substantially outperforming the single-skill baseline at 21.1\%.
This validates that even approximate decomposition provides value for compositional routing.

\paragraph{Metadata matches or outperforms body.}
For oracle decomposition, metadata-only retrieval achieves R@1 = 99.5\% versus 99.0\% for body-aware retrieval.
In the full pipeline, body-aware has a slight edge at the category level ($\catrecall$@1: 36.3\% vs 33.9\%), but the difference is small and comes with significantly higher encoding cost due to longer input sequences.
This suggests that concise skill metadata (name + description) captures the essential discriminative signal, consistent with the intuition that good skill descriptions should be self-contained summaries.

\subsection{Decomposer Comparison}

\begin{table}[t]
\centering
\small
\begin{tabular}{lccc}
\toprule
\textbf{Decomposer} & \textbf{Decomp\%} & $\catrecall$\textbf{@1} & $\chaincat$ \\
\midrule
Qwen2.5-7B & \textbf{0.360} & \textbf{0.339} & \textbf{0.316} \\
Mistral-7B & 0.177 & 0.243 & 0.248 \\
Llama-3.1-8B & 0.000 & 0.216 & 0.210 \\
\midrule
Oracle & 1.000 & 1.000 & 1.000 \\
\bottomrule
\end{tabular}
\caption{Decomposer comparison (all using metadata-only retrieval). Qwen2.5-7B achieves the best decomposition accuracy and downstream routing performance.}
\label{tab:decomposer}
\end{table}

Table~\ref{tab:decomposer} compares the three LLM decomposers.
Qwen2.5-7B significantly outperforms the other models, achieving 36.0\% decomposition accuracy versus 17.7\% for Mistral and 0\% for Llama.

\paragraph{Over-decomposition is the primary failure mode.}
Analysis of decomposition patterns (Figure~\ref{fig:decomp_dist}) reveals that Llama-3.1-8B systematically over-decomposes: its median output is 9 sub-tasks, while the ground truth distribution peaks at 2--3.
No single Llama decomposition matches the ground-truth count, explaining its 0\% decomposition accuracy.
Mistral shows a similar but less severe tendency, while Qwen's distribution most closely aligns with the ground truth.

\subsection{Difficulty Analysis}

\begin{table}[t]
\centering
\small
\begin{tabular}{llccc}
\toprule
\textbf{Method} & \textbf{Diff.} & $\catrecall$\textbf{@1} & $\catrecall$\textbf{@10} \\
\midrule
\multirow{3}{*}{Single-Skill} & Easy & 0.201 & 0.774 \\
& Medium & 0.267 & 0.667 \\
& Hard & 0.146 & 0.675 \\
\midrule
\multirow{3}{*}{Qwen + Meta} & Easy & 0.214 & 0.637 \\
& Medium & 0.367 & 0.658 \\
& Hard & \textbf{0.410} & 0.646 \\
\midrule
\multirow{3}{*}{Qwen + Body} & Easy & 0.252 & 0.607 \\
& Medium & \textbf{0.403} & \textbf{0.728} \\
& Hard & 0.407 & 0.623 \\
\bottomrule
\end{tabular}
\caption{Performance by difficulty level. Decomposition-based routing shows increasing advantage on harder queries.}
\label{tab:difficulty}
\end{table}

Table~\ref{tab:difficulty} breaks down results by difficulty level.
A notable finding is that the decomposition-based pipeline shows its \emph{greatest advantage on harder queries}: for hard queries (4--5 skills), Qwen + Meta achieves $\catrecall$@1 = 41.0\% versus the single-skill baseline at 14.6\%, a relative improvement of 181\%.
For easy queries (2 skills), the improvement is more modest (21.4\% vs 20.1\%).
This pattern suggests that decomposition becomes increasingly important as task complexity grows.

\subsection{Ablation: Constrained Decomposition}

To isolate the effect of decomposition \emph{granularity} (number of steps) from decomposition \emph{quality} (sub-task descriptions), we run an ablation where LLM decompositions are truncated or padded to match the ground-truth step count.

The constrained decomposition achieves $\catrecall$@1 = 33.9\%, identical to the unconstrained pipeline.
This indicates that the primary challenge is not predicting the \emph{number} of steps but generating sub-task descriptions that are \emph{semantically aligned} with the target skills.
Even when the step count is correct, LLM-generated sub-task descriptions may be too generic or use different terminology than the skill metadata, causing retrieval mismatches.

% ============================================================
\section{Analysis}
\label{sec:analysis}

\subsection{Error Analysis}

We manually examine 50 randomly selected failure cases from the Qwen + Meta pipeline to categorize error types:

\begin{itemize}[nosep,leftmargin=*]
  \item \textbf{Over-decomposition} (36\%): The LLM splits a query into more steps than necessary, inserting intermediate steps not in the ground truth (e.g., ``identify necessary patches'' inserted between ``run security audit'' and ``apply patches'').
  \item \textbf{Generic sub-task descriptions} (28\%): The LLM produces descriptions that are too broad (e.g., ``process the data'' instead of ``transform CSV data using pandas''), leading to retrieval of tangentially related skills.
  \item \textbf{Vocabulary mismatch} (22\%): The sub-task description uses different terminology than the skill metadata (e.g., ``create visual reports'' vs. skill named ``matplotlib-chart-generator'').
  \item \textbf{Under-decomposition} (14\%): The LLM merges multiple steps into one, typically for hard queries with 4+ required skills.
\end{itemize}

\subsection{Case Study}

\paragraph{Successful routing (Q0):} 
``\textit{Generate the application code, write tests, and deploy to production}'' is decomposed into three sub-tasks matching the ground truth exactly. The retriever correctly identifies a code generator, a QA tool, and a deployment orchestrator.

\paragraph{Partial success (Q100):}
``\textit{Search for sources, process findings, visualize trends, write report, and export document}'' is correctly decomposed into 5 steps. The retriever correctly identifies skills for search, data processing, and report writing (3/5 categories correct), but selects an exploratory data analysis skill instead of a dedicated visualizer.

\paragraph{Over-decomposition failure (Q200):}
``\textit{Draft a document and localize it for international readers}'' requires 2 skills (writer + translator), but the LLM produces 7 sub-tasks including review, translation, proofreading, and cultural adaptation steps. The first step correctly selects a document writing skill, but subsequent steps diverge from the ground truth.

\subsection{Retrieval Quality with Oracle Decomposition}

The near-perfect oracle performance (R@1 = 99.5\%) deserves further examination.
The 0.5\% of oracle failures occur when multiple skills in the library have very similar metadata, making exact ID matching ambiguous even though the correct \emph{category} of skill is retrieved (category recall remains 100\%).
This validates our design choice of including category-level metrics, which better reflect practical routing quality.

% ============================================================
\section{Discussion}
\label{sec:discussion}

\paragraph{The decomposition gap.}
Our results quantify a ``decomposition gap'': the difference between what is achievable with perfect sub-task descriptions (oracle) and what current LLMs produce.
At $\catrecall$@1, this gap is 66 percentage points (100\% -- 33.9\%).
Closing this gap is the most impactful direction for improving compositional skill routing.
Potential approaches include fine-tuning decomposers on skill-aware data, iterative refinement with retrieval feedback, and leveraging skill metadata as decomposition guidance.

\paragraph{Metadata vs.\ body.}
Our finding that metadata-only retrieval is competitive with body-aware retrieval has practical implications: it means that skill authors should invest in writing clear, descriptive metadata (names and descriptions) rather than relying on the full specification body to be the primary retrieval signal.
This finding contrasts with \citet{xu2025skillrouter}, who report that body content is the ``decisive'' signal for single-skill routing.
We hypothesize that the compositional setting, where queries are decomposed into specific sub-tasks, produces more targeted retrieval queries that align naturally with concise metadata.

\paragraph{Scalability.}
Our FAISS-based retrieval scales to thousands of skills with sub-second latency per query.
The primary computational cost is LLM-based decomposition, which at approximately 0.8 seconds per query (sequential generation on A100) is practical for interactive use.
Batch processing reduces amortized cost further.

\paragraph{Limitations.}
Our benchmark uses template-based query generation, which may not fully capture the diversity of real-world compositional queries.
The sequential DAG assumption, while covering most cases in our benchmark, does not model parallel skill execution.
We evaluate with 7--8B parameter LLMs; larger models may achieve better decomposition quality.
Finally, our evaluation is retrieval-focused and does not measure end-to-end task completion with actual skill execution.

% ============================================================
\section{Conclusion}
\label{sec:conclusion}

We have formalized the Compositional Skill Routing problem and presented \system{}, a decompose-retrieve-compose framework for routing complex queries to multiple skills.
Through \benchmark{}, a new benchmark of 300 compositional queries over 2{,}595 curated skills, we demonstrate that task decomposition quality is the primary bottleneck in compositional routing: oracle decomposition achieves near-perfect retrieval, while the best 7B LLM decomposer reaches 33.9\% category recall.
Our analysis identifies over-decomposition and vocabulary mismatch as the dominant failure modes, providing clear directions for future work on decomposition-aware skill routing.

% ============================================================
% References
\bibliography{references}

% ============================================================
\appendix

\section{Category Taxonomy}
\label{app:categories}

The 17 functional categories in \benchmark{} are: 
code\_generator, api\_client, test\_writer, data\_processor, code\_analyzer, writer, deployment, security, visualizer, search, refactorer, file\_converter, translator, database, monitoring, designer, git\_workflow.

\section{Decomposition Distribution}
\label{app:decomp}

Figure~\ref{fig:decomp_dist} shows the distribution of predicted sub-task counts for each LLM decomposer compared to the ground truth.

\begin{figure}[h]
\centering
% Placeholder for actual figure
\fbox{\parbox{0.9\columnwidth}{\centering\vspace{2em}
\small Distribution of sub-task counts:\\
GT: peak at 2--3 | Qwen: peak at 3--4\\
Mistral: peak at 3--5 | Llama: peak at 8--11
\vspace{2em}}}
\caption{Distribution of predicted sub-task counts by decomposer. Llama-3.1-8B systematically over-decomposes, with a median of 9 sub-tasks versus 3 in the ground truth.}
\label{fig:decomp_dist}
\end{figure}

\section{Top-$k$ Sensitivity}
\label{app:topk}

With oracle decomposition, R@1 = 99.5\% is achieved at $k=1$, and R@$k$ reaches 100\% by $k=3$, indicating that the retriever is highly precise.

\end{document}
